\section{Deployment}
\label{sec:deployment}

The CrowdFund DApp supports two deployment modes: a local development environment using Foundry's Anvil and a production deployment on Vercel.

\subsection{Local Deployment (Anvil)}

\subsubsection{Step 1: Start Local Blockchain}

Anvil provides a local Ethereum node with 10 pre-funded accounts, instant block mining, and deterministic private keys:

\begin{lstlisting}[language={}, caption={Starting Anvil}]
anvil
\end{lstlisting}

Anvil starts on \texttt{http://127.0.0.1:8545} with Chain ID \texttt{31337}.

\subsubsection{Step 2: Deploy Contract}

The deployment script (\texttt{script/deploy.s.sol}) uses Foundry's \texttt{Script} base contract:

\begin{lstlisting}[language=Solidity, caption={Deployment Script}]
contract DeployCrowdFunding is Script {
    function run() external returns (CrowdFunding crowdFunding) {
        vm.startBroadcast();
        crowdFunding = new CrowdFunding();
        vm.stopBroadcast();
    }
}
\end{lstlisting}

Execution:
\begin{lstlisting}[language={}, caption={Deploy Command}]
forge script script/deploy.s.sol:DeployCrowdFunding \
    --rpc-url http://127.0.0.1:8545 \
    --private-key <ANVIL_PRIVATE_KEY> \
    --broadcast
\end{lstlisting}

An automated shell script (\texttt{deploy.sh}) loads environment variables from \texttt{.env} and executes the deployment, avoiding manual entry of sensitive keys.

\subsubsection{Step 3: Configure MetaMask}

MetaMask is configured to connect to the local Anvil node:

\begin{table}[H]
\centering
\caption{MetaMask Network Configuration}
\begin{tabular}{l l}
\toprule
\textbf{Parameter} & \textbf{Value} \\
\midrule
Network Name    & Anvil Local \\
RPC URL         & \texttt{http://127.0.0.1:8545} \\
Chain ID        & 31337 \\
Currency Symbol & ETH \\
\bottomrule
\end{tabular}
\end{table}

\subsubsection{Step 4: Start Frontend}

\begin{lstlisting}[language={}, caption={Frontend Development Server}]
cd frontend
npm run dev
\end{lstlisting}

Vite starts a development server with hot module replacement at \texttt{http://localhost:5173}.

\subsection{Production Deployment (Vercel)}

For production, the frontend and serverless API functions are deployed to Vercel:

\subsubsection{Architecture}

\begin{itemize}[leftmargin=1.5em]
    \item \textbf{Static frontend:} Vite builds the React application into the \texttt{dist/} directory, which is served as static files.
    \item \textbf{Serverless API:} The \texttt{api/} directory contains serverless functions (\texttt{upload.js}, \texttt{ipfs/[id].js}) that replace the Express server in production.
    \item \textbf{Vercel Blob Storage:} In production, campaign metadata is stored in Vercel Blob Storage instead of the local filesystem, ensuring persistence across serverless function invocations.
\end{itemize}

\subsubsection{Routing Configuration}

\begin{lstlisting}[language={}, caption={vercel.json}]
{
  "buildCommand": "vite build",
  "outputDirectory": "dist",
  "framework": "vite",
  "routes": [
    { "src": "/api/upload", "dest": "/api/upload.js" },
    { "src": "/api/ipfs/(.+)", "dest": "/api/ipfs/[id].js?id=$1" },
    { "handle": "filesystem" },
    { "src": "/(.*)", "dest": "/index.html" }
  ]
}
\end{lstlisting}

\subsubsection{Environment Variable Security}

All sensitive credentials are managed via Vercel's environment variable system, never committed to version control:

\begin{itemize}[leftmargin=1.5em]
    \item \texttt{PINATA\_API\_KEY} and \texttt{PINATA\_API\_SECRET}: IPFS pinning service credentials.
    \item \texttt{BLOB\_READ\_WRITE\_TOKEN}: Vercel Blob Storage access token.
    \item Private keys are \emph{never} stored in the production environment; all transactions are signed client-side via MetaMask.
\end{itemize}

\subsection{Build and Test Automation}

A \texttt{build\_and\_test.sh} script automates the full CI pipeline:

\begin{lstlisting}[language={}, caption={CI Script}]
#!/bin/bash
set -e
forge build          # Compile contracts
forge test -vvv      # Run tests with verbose output
forge test --gas-report  # Generate gas report
\end{lstlisting}
